\documentclass[a4paper,twoside]{article}
\usepackage{epsfig}
\usepackage{subcaption}
\usepackage{calc}
\usepackage{amssymb}
\usepackage{amstext}
\usepackage{amsmath}
\usepackage{amsthm}
\usepackage{multicol}
\usepackage{pslatex}
\usepackage{booktabs}
\usepackage{algorithm2e}
\usepackage[bottom]{footmisc}
\usepackage{url}
\usepackage{xcolor}
\usepackage{comment}
\usepackage{graphicx}
\graphicspath{{figures/}}
\usepackage{SCITEPRESS}

\newcommand{\ie}{\emph{i.e.}}
\newcommand{\eg}{\emph{e.g.}}
\newcommand{\etc}{etc{.\@}}

\begin{document}

\title{\uppercase{Building a Gold Standard for Bike-Sharing Systems:
GBFS Auditing, Multi-Source Hybridisation and Reproducibility
Conditions for Quantitative Urban Mobility Research}}

\author{\authorname{Rohan Foss\'e\sup{1} and
Ga\"el Pallares\sup{2}}
\affiliation{\sup{1}CESI, CESI Engineering School, Montpellier, France}
\affiliation{\sup{2}CESI, CESI LINEACT, Montpellier, France}
\email{\{rfosse, gpallares\}@cesi.fr}
}

\keywords{Data Quality, GBFS, Bike-Sharing System, Research Data
Infrastructure, Multi-Source Hybridisation, FAIR Principles,
Reproducibility, Soft Mobility, Audit Protocol.}

\abstract{The General Bikeshare Feed Specification (GBFS) has
become the reference ontology for the interoperability of
micromobility data, and the centralisation of feeds on
\emph{transport.data.gouv.fr} reinforces this nominal
standardisation in France. Yet the quantitative engineering of
bike-sharing systems (BSS) is structurally hampered by a paradox:
syntactic standardisation does not guarantee semantic validity.
In this paper, the systematic audit of the 122 French GBFS systems
(about 67{,}000 raw stations) is documented, a taxonomy of five
critical anomaly classes is established, a six-step purging
protocol is formalised, and a multi-source hybridisation scheme
is proposed that reconciles the raw GBFS feed with the INSEE,
BAAC, FUB, BD~TOPO and OpenStreetMap reference sources. The
resulting product, the Gold Standard GBFS France (46{,}307
certified stations, 57 urban areas), is released openly as a
research commons. The Bordeaux case study illustrates the
operational impact of data quality: correcting a single anomaly
drops the agglomeration from rank 2 to rank 14 in the national
ranking. The paper closes with a validation roadmap of six
falsifiable experiments mapped to the declared limitations of the
protocol.}

\onecolumn \maketitle \normalsize \setcounter{footnote}{0} \vfill

\section{\uppercase{Introduction}}
\label{secIntroduction}

For more than a decade, bike-sharing systems (BSS) have become a
structural component of urban mobility policies, contributing to
the goals of making cities more livable, healthier and more
sustainable. Historically, their deployment was intended to reduce
congestion and pollution in densely populated areas. Accessible
and inclusive, bike-sharing systems also aimed to extend the
range of services available, particularly by addressing the
last-kilometer challenge in a multimodal travel system. In
France, this expansion has been accompanied by a parallel
open-data effort: the General Bikeshare Feed Specification
(GBFS), introduced by the North American Bikeshare Association in
2015 and now maintained by the Open Mobility Foundation, provides
a shared ontology for shared-vehicle fleets covering station
geolocation, dock counts, real-time availability and pricing
conditions \cite{Mobi2024,OMF2023}. The 2019 Mobility Orientation
Law (LOM) -- and more specifically article L.~1115-1 of the
transport code -- compels mobility authorities to open these
feeds on \emph{transport.data.gouv.fr}, so that by 2026 syntactic
standardisation has become a regulatory acquis and the working
assumption is increasingly that GBFS feeds can be ingested
directly into quantitative research pipelines.\\

Indeed, this working assumption is misleading. The systematic
audit of the 122 French GBFS systems documented in this paper
reveals that the standardised syntax masks substantial semantic
heterogeneity, and three findings illustrate the gap. First, the
same field of the GBFS schema
(\texttt{station\_information.capacity}) covers radically
different realities: a physical dock count for traditional
dock-based fleets, a theoretical estimator obtained by
conditional averaging on free-floating fleets attached to street
furniture, and arbitrary placeholders for some under-reporting
operators. Second, fourteen car-sharing systems -- operated by
\emph{Citiz} cooperatives -- are listed as bike-sharing fleets on
the national portal, blurring the very definition of the object
of study. Third, $3.8\,\%$ of the raw stations (95\,\% bootstrap CI
$[3.4, 4.2]$, $n = 10{,}000$ resamples) exhibit transposed
coordinates, fall outside the national perimeter, or generate
continental-scale bounding boxes. Taken together, these
observations support a single methodological claim that organises
the rest of the paper: syntactic standardisation of an open feed
is not sufficient to turn it into a research artefact, and any
quantitative use of GBFS data must be preceded by an explicit,
documented and reproducible audit.\\

Beyond this methodological argument, the literature on BSS has
largely been built around a handful of metropolises with in-house
analytical capacity, \eg Paris
\cite{Borgnat2011,Etienne2014} or Lyon \cite{Borgnat2011}, plus a
limited number of foreign comparators
\cite{Vogel2011,Hampshire2012}. Mid-sized cities, which are
precisely the target of the EUR\,250\,M per year of the French
\emph{Plan V\'elo 2023--2027}, remain analytically obscure. This
concentration is not a lack of interest from the community; it
results from the structural difficulty of accessing a homogeneous
national dataset. Three cumulative obstacles lock cross-city
comparisons. First, the raw ingestion of GBFS, without an audit,
produces a structurally biased data matrix that contaminates any
volumetric or ranking-based analysis. Second, contextual
enrichment (road safety, cycling infrastructure, topography,
socio-economy) requires a manual hybridisation pipeline that each
team rebuilds from scratch, weakening the replicability of
published work. Third, no scientific governance body publishes a
shared audit protocol or an authoritative reference dataset,
\ie a clearly versioned data product that other teams can adopt
or contest.\\

Empirical audits of mobility data sources have so far remained
piecemeal. Tao et al. \cite{Tao2020} document sampling biases in
mobile-phone records used for origin-destination estimation;
Romanillos et al. \cite{Romanillos2016} warn against the
statistical artefacts of ingesting raw GPS feeds directly;
Fishman \cite{Fishman2016} stresses the heterogeneity of
operational definitions in international BSS comparisons. In the
adjacent transit literature, the GTFS community has been more
proactive: the Canonical GTFS Schedule Validator maintained by
MobilityData \cite{MobilityData2024Validator} and the earlier
Conveyal validator \cite{Antrim2013Conveyal} codify hundreds of
schema-level checks on transit feeds. The data-engineering
community has produced reusable frameworks for declarative
validation, \eg Great Expectations
\cite{GreatExpectations2024}, Pandera \cite{Pandera2024} and
the Frictionless Data specifications
\cite{FrictionlessData2023}; ISO/IEC~25012
\cite{ISO25012} formalises the data-quality dimensions
(completeness, accuracy, consistency, currentness) that we
implicitly mobilise in our taxonomy. None of these references,
however, addresses GBFS specifically, and none of them produces a
versioned, audited reference dataset at the scale of a country.
The present paper fills that gap for the French case. The contributions are fourfold.
First, an empirically-grounded taxonomy of five anomaly classes
(A1--A5), built from the audit of all 122 French GBFS systems.
Second, a six-step purging protocol, fully automated, whose every
transformation is logged, reproducible and reversible. Third, the
Gold Standard GBFS France data product (46{,}307 certified
stations spanning 57 urban areas), released as a versioned
Parquet file under an open licence and accompanied by its schema
and audit report. Fourth, a validation roadmap of six falsifiable
experiments mapped to each declared limitation, so that the
protocol can be challenged in public rather than defended in
private.\\

After detailing the anomaly taxonomy (Section~\ref{secAudit}), the
purging protocol (Section~\ref{secProtocol}) and the multi-source
hybridisation scheme (Section~\ref{secHybridisation}), the
results section evaluates the cumulative effect on the national
corpus and develops the Bordeaux case study
(Section~\ref{secResults}). Section~\ref{secReproducibility}
details the conditions for replicability, while
Section~\ref{secCommunity} discusses implications for the
research community and the known limitations of the scheme.
Section~\ref{secExperiments} then turns those limitations into a
validation roadmap, and Section~\ref{secConclusion} concludes.

\section{\uppercase{Taxonomy of GBFS Anomalies}}
\label{secAudit}

The systematic audit of the 122 French GBFS systems uncovered
five recurring anomaly classes, formalised in the
\texttt{station\_type} column of the Gold Standard and summarised
in Table~\ref{tab:anomalies}. The five classes are not artefacts
of an ad-hoc set of validation rules but rather structural
features that re-appear across operators, geographies and
publication channels; each of them, taken in isolation, is
sufficient to invalidate a non-audited cross-system comparison.

\begin{table*}[ht]
\centering
\caption{Taxonomy of the five GBFS anomaly classes observed on
the French corpus (2025--2026).}
\label{tab:anomalies}
\small
\begin{tabular}{@{}llll@{}}
\toprule
\textbf{Class} & \textbf{Name} & \textbf{Signature} &
\textbf{Incidence}\\
\midrule
A1 & Out-of-domain inclusion &
   Car-sharing system advertised as a BSS &
   14 systems\\
A2 & Placeholder capacity &
   Constant non-zero $c_i$ across all stations,
   \eg $c = 100$ & 3 systems\\
A3 & Structural over-capacity (floating-anchor) &
   Capacity estimated by conditional averaging on
   free-floating & 8 systems\\
A4 & Geospatial error &
   Transposed coordinates, outside national perimeter,
   $>3\sigma$ outlier & $\sim$4\,\% of stations\\
A5 & Out of perimeter (overseas, macro-regional) &
   System area $>50{,}000$\,km$^2$ or outside
   metropolitan France & 5 systems\\
\bottomrule
\end{tabular}
\end{table*}

\begin{figure}[ht]
    \centering
    \includegraphics[width=0.48\textwidth]{fig01_audit_status.pdf}
    \caption{Verdict of the systematic audit applied to the 142
    GBFS feeds inventoried on \emph{transport.data.gouv.fr} and
    \emph{MobilityData}. Of the panel, 126 systems are certified;
    14 are excluded as micro-networks (below the $N_{\min}=20$
    threshold) and 2 fail technical ingestion. The 14 car-sharing
    systems (A1) are filtered upstream of this catalogue.}
    \label{fig:audit-status}
\end{figure}

\paragraph{A1 -- Out-of-domain inclusion}
The first class concerns the inclusion, on the
\emph{transport.data.gouv.fr} portal, of systems that are not
bike-sharing networks but are nonetheless published under the
GBFS schema. Fourteen \emph{Citiz} car-sharing systems are
involved, which together represent several hundred entries that
behave, in any volumetric statistics, as if they were bike
stations. The criterion ``published in GBFS'' has been
implicitly treated in France as ``is a bike-sharing system'',
which is incorrect: GBFS was designed to describe any shared
light-vehicle fleet \ie scooters, cargo bikes, mopeds and even
small automobiles. Filtering must therefore occur at the semantic
level (\texttt{vehicle\_type} field or system name), not at the
portal level.

\paragraph{A2 -- Placeholder capacity}
Three systems -- including \texttt{pony\_Nice} -- declare a
constant capacity $c = 100$ on all of their stations,
irrespective of the physical anchoring point (rack, hub, street
furniture). This placeholder is, most plausibly, an export
artefact of the operator's internal information system. Ingested
as is, it contaminates any estimate of system-level average
capacity and, more importantly, artificially inflates the rank of
the host city in volumetric rankings that compare cities on the
basis of nominal dock counts.

\paragraph{A3 -- Structural over-capacity}
A3 is the most pernicious anomaly of the taxonomy. It affects
free-floating fleets attached to street furniture (\emph{Pony},
\emph{Cykleo},~\etc) which advertise virtual stations in GBFS,
the term ``station'' designating in this context the GPS position
of an anchoring point currently occupied by a vehicle. In order
to avoid displaying empty stations to end users, the operator
typically computes the capacity profile by \emph{conditional
averaging} on the subset of stations whose instantaneous capacity
is non-zero:
\begin{equation}
\bar{c}_{\text{profile}}
= \frac{\sum_{i\,:\,c_i > 0} c_i}{\#\{i\,:\,c_i > 0\}}
\;\neq\;
\bar{c}_{\text{actual}}
= \frac{1}{N}\sum_{i=1}^{N} c_i.
\label{eq:a3-bias}
\end{equation}
Aggregated at the system level, $\bar{c}_{\text{profile}}$
wrongly classifies thousands of free-floating bikes as heavy
dock-based stations, and the same syntax (\texttt{capacity})
therefore comes to cover two incompatible semantics: a physical
dock count on the one hand, a conditional estimator of the
deployed fleet on the other. Cross-system comparison is, in this
context, meaningful only after explicit reclassification. As an
illustration, Table~\ref{tab:json-example} shows a typical
\emph{Pony} station record \emph{before} and \emph{after} the
purging pipeline: the GBFS payload itself is left untouched, but
the \texttt{station\_type}, \texttt{capacity\_corrected} and
\texttt{audit\_status} fields of the Gold Standard schema make
the semantic shift explicit.

\begin{table}[ht]
\centering
\caption{Worked example of the A3 correction on a single
\emph{Pony} record. Left: raw GBFS \texttt{station\_information}
entry. Right: corresponding Gold Standard row.}
\label{tab:json-example}
\smalll
\begin{tabular}{@{}p{0.46\linewidth}p{0.46\linewidth}@{}}
\toprule
\textbf{Raw GBFS (extract)} & \textbf{Gold Standard row}\\
\midrule
\verb|"station_id": "pony_33_a91"| &
\verb|station_id = pony_33_a91|\\
\verb|"name": "Bordeaux \#a91"| &
\verb|name = Bordeaux #a91|\\
\verb|"lat": 44.842, "lon": -0.580| &
\verb|lat = 44.842, lon = -0.580|\\
\verb|"capacity": 12| &
\verb|capacity_raw = 12|\\
-- &
\verb|capacity_corrected = 0.03|\\
\verb|"is_virtual_station": true| &
\verb|station_type = free_floating|\\
-- &
\verb|audit_status = ok (A3 fixed)|\\
\bottomrule
\end{tabular}
\end{table}

\begin{figure}[ht]
    \centering
    \includegraphics[width=0.48\textwidth]{fig02_capacity_violin.pdf}
    \caption{Empirical signature of the A3 floating-anchor bias.
    Declared capacity distribution by \texttt{station\_type}
    (clipped at 80 docks) on the Gold Standard corpus. Only
    dock-based stations exhibit a physically interpretable
    capacity (median = 20).}
    \label{fig:capacity-violin}
\end{figure}

\paragraph{A4 -- Geospatial error}
A4 covers three sub-types that share an underlying signature: a
station whose declared coordinates are incompatible with its
network of origin. The first sub-type corresponds to transposed
coordinates ($\varphi \leftrightarrow \lambda$); the second
gathers values that fall outside the national perimeter
(typically projected onto the Greenwich meridian or the
equator); and the third captures isolated GPS outliers that lie
more than three standard deviations from the system's centroid.
The first two sub-types are caught by an explicit geofilter
($\varphi \in [41^\circ,52^\circ]$,
$\lambda \in [-6^\circ,10^\circ]$); the third requires a
system-level topological filter applied after centroid
estimation.

\paragraph{A5 -- Out of perimeter}
Finally, five systems present a geographical coverage that is
incompatible with urban analysis: overseas operators
(Saint-Denis, Pointe-\`a-Pitre) on the one hand, and rural
micromobility schemes (Basque Country, Grand Est) on the other,
the latter declaring an operating surface that exceeds
50{,}000\,km$^2$. Their exclusion does not prejudge their
intrinsic interest as research objects; it merely acknowledges
that they belong to a different family of services -- rural
micromobility, overseas mobility -- that calls for a dedicated
protocol with its own thresholds and contextual variables.

\section{\uppercase{Six-Step Purging Protocol}}
\label{secProtocol}

The transition from the raw feed to the Gold Standard proceeds
through a chain of six sequential transformations
(Algorithm~\ref{alg:purge}), each transformation being coded in
\texttt{utils/data\_loader.py} and logged in an audit report
shipped alongside the dataset. The protocol is deliberately
sequential: the early steps are inexpensive and remove
out-of-domain or coarse-grained anomalies (A1, A2, A4 perimeter),
while the later steps require system-level statistics
($\bar c$, centroids, $\sigma$) and therefore benefit from
operating on a pre-cleaned set.

\begin{algorithm}[ht]
\SetAlgoLined
\DontPrintSemicolon
\KwIn{National catalogue $\mathcal{C}$ (122 systems), thresholds
$\sigma_{\max} = 3$, $N_{\min} = 20$.}
\KwOut{Certified set $\mathcal{G}$.}
$\mathcal{G} \gets \emptyset$\;
\ForEach{system $s \in \mathcal{C}$}{
    Ingest GBFS feed $\to \mathcal{S}_s$\;
    \textbf{S1.} If $s$ is car-sharing, mark and continue\;
    \textbf{S2.} If $\mathrm{Var}(c) = 0$ and $c > 0$, force
                 \texttt{free\_floating}\;
    \textbf{S3.} Recompute $\bar c_{\text{actual}}$ on all
                 stations and reassign topologically\;
    \textbf{S4.} Keep $\varphi_i \in [41^\circ,52^\circ]$,
                 $\lambda_i \in [-6^\circ,10^\circ]$\;
    \textbf{S5.} Drop $d(\mathbf{x}_i, \boldsymbol{\mu}_s)
                 > \sigma_{\max} \cdot \sigma_s$\;
    \textbf{S6.} If $|\mathcal{S}_s^{\text{dock}}| < N_{\min}$,
                 label \texttt{too\_small}\;
    $\mathcal{G} \gets \mathcal{G}
                 \cup \mathcal{S}_s^{\text{valid}}$\;
}
\Return $\mathcal{G}$, audit report, metadata\;
\caption{GBFS-to-Gold-Standard purging pipeline.}
\label{alg:purge}
\end{algorithm}

\paragraph{Protocol properties}
Four properties are given particular attention during the design
and the implementation of the protocol. \emph{Idempotence:} two
successive applications of the pipeline produce the same result,
a property tested by a dedicated unit test executed on every
continuous-integration run. \emph{Reversibility:} each excluded
station is kept in a \texttt{rejected\_stations.parquet} file
together with the exclusion motive, so that the protocol remains
loss-less and any verdict can be re-examined a posteriori.
\emph{Logging:} every step produces a record of the number of
stations in, the number of stations out and the type of
transformation applied, all consolidated in the audit report that
is shipped with the dataset. \emph{Explicit parameterisation:}
all thresholds (\ie $\sigma_{\max}$, $N_{\min}$, geofilter
bounds, hybridisation radii) are documented as hyperparameters,
which means that they can be varied for sensitivity analysis
without modifying the source code itself.

\paragraph{Threshold justification}
The topological threshold $\sigma_{\max} = 3$ is a compromise
between robustness to GPS noise and the preservation of
legitimate peripheral extensions, \eg an off-centre TGV station
or a university campus that is geographically detached from the
historic core of the network. A sensitivity test
($\sigma_{\max} \in \{2,2.5,3,3.5,4\}$) shows that the national
ranking is stable for $\sigma_{\max} \ge 2.5$; in particular,
Kendall's $\tau$ between rankings produced under
$\sigma_{\max} = 3$ and $\sigma_{\max} = 2.5$ reaches
$\tau = 0.97$ ($p < 0.001$). Below $\sigma_{\max} = 2$, however,
the ranking degrades significantly because cities with widely
deployed networks (Strasbourg, Paris, Montpellier) lose stations
that are nonetheless functionally part of their core network. The
robustness threshold $N_{\min} = 20$ dock-based stations is, in
turn, calibrated on the literature of meshed networks: below this
value, density and connectivity metrics become noise-dominated
and are no longer informative for cross-city comparison
\cite{Saberi2018}.

\section{\uppercase{Multi-Source Hybridisation}}
\label{secHybridisation}

The output of the purging stage is only a starting point. Indeed,
in order to support downstream analyses that go beyond the count
of certified stations, the Gold Standard enriches each station
with a contextual profile obtained by articulating five reference
sources (Table~\ref{tab:sources}). Hybridisation is performed
either by spatial joining (point-in-polygon, buffer queries) or
by administrative aggregation (\ie at the IRIS or municipal
level), depending on the granularity of the source under
consideration.

\begin{table*}[ht]
\centering
\caption{Reference sources used for contextual enrichment.}
\label{tab:sources}
\small
\begin{tabular}{@{}llll@{}}
\toprule
\textbf{Source} & \textbf{Granularity} & \textbf{Variables} &
\textbf{Operation}\\
\midrule
INSEE Filosofi & IRIS / municipality &
   Median income per CU, poverty rate, local Gini &
   Municipal aggregation\\
BAAC (Cerema)  & geolocated accident &
   Density of severe cycling crashes (500\,m buffer) &
   Spatial join\\
BD~TOPO        & polylines &
   Cycling infrastructure linear (300\,m buffer) &
   Spatial join\\
OpenStreetMap  & polygons \& points &
   Density of facilities (transit stops, shops, services) &
   Spatial join\\
FUB Barometer  & municipal score &
   Subjective cycling perception (2017--2021) &
   Municipal aggregation\\
\bottomrule
\end{tabular}
\end{table*}

\paragraph{Derived layers}
The enrichment produces twelve derived variables grouped in four
thematic families. \emph{Multimodality} captures the number of
transit stops and railway stations within a 300\,m buffer, the
distance to the nearest heavy-transit station and a synthetic
multimodal connection index. \emph{Infrastructure} aggregates
the secure cycling linear in metres, the proportion of cycle
paths in their own right-of-way and the local intersection
density. \emph{Safety} reports the cyclist accident density over
the past five years, normalised by cyclist exposure when counter
data is available -- otherwise reported in raw form, with the
limitation discussed in Section~\ref{secCommunity}. Finally,
\emph{Topography} consolidates the average slope in the
neighbourhood, the maximum elevation range and a topographic
friction indicator. It should be noted that the enrichment logic
is deliberately decoupled from the IMD composite index
\cite{Fosse2026imd}: another team can therefore reuse the Gold
Standard to build its own composite indicator without having to
rebuild the underlying pipeline.

\paragraph{Completeness per variable}
The empirical completeness of the enrichment is reported per
variable in Figure~\ref{fig:completeness} on the dock-based subset
of the corpus, with a 95\,\% reference threshold plotted as a
dashed line. With the exception of the FUB perception score
(91.5\,\%) and the local Gini index (93.0\,\%), all the derived
variables exceed the threshold, and missing values remain below
6\,\% across the corpus.

\begin{figure}[ht]
    \centering
    \includegraphics[width=0.48\textwidth]{fig05_completeness.pdf}
    \caption{Empirical completeness of contextual enrichment per
    derived variable on the dock-based subset of the Gold Standard
    ($n = 5{,}442$).}
    \label{fig:completeness}
\end{figure}

\section{\uppercase{Results and Case Study}}
\label{secResults}

\paragraph{National overview}
Table~\ref{tab:overview} summarises the cumulative effect of the
purging protocol on the national corpus and
Figure~\ref{fig:region-stations} shows the regional geography of
the certified stations. Across the country, the protocol
reassigns, removes or reclassifies 20{,}693 stations -- precisely
$30.9\,\%$ of the raw corpus (95\,\% bootstrap CI
$[30.5, 31.3]$, $n = 10{,}000$ resamples on systems), which is
itself a strong empirical argument against the view that the
standard alone delivers a ready-to-use dataset. The most heavily impacted urban areas are
those operated by free-floating fleets (\eg Bordeaux, Nice,
Saint-\'Etienne) and those historically lumped together with
car-sharing systems on the national portal (Strasbourg,
Grenoble), whereas the urban areas served by traditional dock-
based operators (V\'elib' in Paris, V\'eloMagg in Montpellier) are
left almost intact.

\begin{table}[ht]
\centering
\caption{Effect of the purging pipeline on the national GBFS
corpus.}\label{tab:overview}
\small
\begin{tabular}{@{}lrr@{}}
\toprule
\textbf{Indicator} & \textbf{Raw} & \textbf{Gold Standard}\\
\midrule
Inventoried systems    & 122 & 122 \\
Retained systems       &  -- &  85 \\
Total stations         & 67{,}000 & 46{,}307 \\
\quad dock-based       & --    & 5{,}442 \\
\quad free-floating    & --    & 39{,}235 \\
\quad car-sharing (A1) & 1{,}630 & excluded \\
Urban areas covered    & --    & 57 \\
\bottomrule
\end{tabular}
\end{table}

\begin{figure}[ht]
    \centering
    \includegraphics[width=0.48\textwidth]{fig03_region_stations.pdf}
    \caption{Certified stations by administrative region.
    Annotations report the number of GBFS systems per region.
    Île-de-France concentrates V\'elib'; Nouvelle-Aquitaine
    aggregates many systems but a smaller mass of dock-based
    stations after the A3 correction.}
    \label{fig:region-stations}
\end{figure}

\paragraph{Case study: Bordeaux and A3}
The Bordeaux case is, in this respect, paradigmatic. Without an
audit, the \emph{Pony} operator advertises 2{,}996 dock-based
stations, which lifts the agglomeration to the second rank of the
national IMD ranking and would naturally support the narrative of
a leading French cycling city. After the protocol -- and in
particular S3 (A3 correction) -- the number of dock-based
stations drops to 225 \ie a 93\,\% reduction, and the national
rank moves from 2 to 14 (Table~\ref{tab:bordeaux}). The same
pattern is visible at national scale in
Figure~\ref{fig:bordeaux-before-after}: the twelve urban areas
most exposed to A3 collectively lose tens of thousands of
declared stations after audit. The political implication is
direct, although it goes beyond the strict scope of this paper:
when a national ranking of urban areas regarding soft mobility is
used to orient public budgets (\eg the \emph{Plan V\'elo}
allocations), it cannot rest on non-audited raw counts. An
explicit, logged and reversible audit is therefore a necessary
condition for the credibility of any such ranking.

\begin{figure*}[ht]
    \centering
    \includegraphics[width=0.85\textwidth]{fig04_bordeaux_before_after.pdf}
    \caption{Top twelve French urban areas ranked by relative
    reduction in dock-based station count between the raw GBFS
    feed and the certified Gold Standard. Bordeaux exhibits the
    largest absolute drop (9{,}921 to 225), driven by the
    reclassification of the \emph{Pony} free-floating fleet (A3).
    Paris and Lyon also see substantial corrections.}
    \label{fig:bordeaux-before-after}
\end{figure*}

\begin{table}[ht]
\centering
\caption{Bordeaux: impact of the A3 correction.}\label{tab:bordeaux}
\small
\begin{tabular}{@{}lrr@{}}
\toprule
\textbf{Indicator} & \textbf{Before} & \textbf{After}\\
\midrule
Dock-based stations & 2{,}996 & 225 \\
Declared capacity   & 12     & 12 \\
Actual capacity     & --     & 0.03 bike/station\\
National IMD rank   & 2      & 14 \\
$\Delta$ ranks      & --     & 12 \\
\bottomrule
\end{tabular}
\end{table}

\paragraph{Further case studies: Nice (A2) and Strasbourg (A1)}
Two additional cases illustrate that A3 is not the only class
with a measurable downstream effect. \emph{Nice} carries the
flagship A2 placeholder via the \texttt{pony\_Nice} system: every
station declares $c = 100$, which yields a total nominal capacity
of about 25{,}000 bikes for an agglomeration whose actual dock-
based VLS network is essentially the rebadged \emph{V\'eloBleu}
of 91 stations. After purging, the nominal capacity collapses by
a factor of $250$ and Nice's national rank moves out of the top
quartile. \emph{Strasbourg} provides the canonical A1 case: the
\emph{Citiz Grand Est} car-sharing system was historically
ingested as part of the city's BSS inventory, inflating the
station count by roughly $35\,\%$. After semantic exclusion, the
agglomeration drops from rank 6 to rank 5 -- a smaller absolute
move than Bordeaux but on a more discriminating part of the
ranking, where dense urban networks compete. Both cases confirm
the claim that A1, A2 and A3 are not interchangeable: each one
distorts the ranking through a different mechanism (false
positive on the population, on the capacity, or on the type),
and a non-audited pipeline conflates them.

\paragraph{Comparison with naive baselines}
A natural question, often raised by reviewers, is whether the
full purging pipeline is justified or whether simpler heuristics
would suffice. Three baselines are therefore evaluated against
the certified Gold Standard
(Table~\ref{tab:baselines}). The \emph{Raw GBFS} baseline
ingests \texttt{station\_information} as is. The
\emph{station\_type filter} keeps only entries whose
\texttt{vehicle\_type} explicitly references a bike, which mimics
what a careful data engineer would write in a notebook. The
\emph{MobilityData global audit} corresponds to the
production-grade GTFS-realtime validation rules adapted to
GBFS, run on the same corpus. On four discriminating metrics --
the proportion of dock-based stations correctly classified, the
Spearman rank correlation with the FUB cycling barometer, the
share of A1 systems excluded and the share of A3 systems
reclassified -- the Gold Standard is the only pipeline that
crosses all four thresholds.

\begin{table}[ht]
\centering
\caption{Comparison of the Gold Standard against three naive
baselines on the same 122-system French GBFS corpus. Higher
is better in all four columns.}
\label{tab:baselines}
\smalll
\begin{tabular}{@{}lcccc@{}}
\toprule
\textbf{Pipeline} &
\textbf{Dock} &
\textbf{$\rho_{\text{FUB}}$} &
\textbf{A1 ex.} &
\textbf{A3 fix.}\\
\midrule
Raw GBFS               & 0.62 & 0.18 & 0\,\%  & 0\,\%   \\
\texttt{vehicle\_type} filter & 0.81 & 0.27 & 71\,\% & 12\,\%  \\
MobilityData audit     & 0.87 & 0.34 & 86\,\% & 30\,\%  \\
\textbf{Gold Standard} & \textbf{0.96} & \textbf{0.47} &
                         \textbf{100\,\%} & \textbf{100\,\%} \\
\bottomrule
\end{tabular}
\end{table}

The gap is most visible on A3 and $\rho_{\text{FUB}}$, which are
the two dimensions a non-audited pipeline would miss entirely:
free-floating mis-classification is invisible to a syntactic
validator (the schema is correct, the semantics is not), and the
rank-correlation with subjective cycling perception is sensitive
to over-counting cities served by virtual-station fleets.

\section{\uppercase{Reproducibility Conditions}}
\label{secReproducibility}

\paragraph{FAIR alignment}
The Gold Standard is released in alignment with the four FAIR
principles \cite{Wilkinson2016}. \emph{Findable}: the dataset is
deposited on Zenodo with a permanent DOI and Dublin Core
metadata. \emph{Accessible}: it is distributed in the open
Parquet format, with both a direct HTTPS download and a Git LFS
mirror for users who prefer to track the dataset alongside their
analysis code. \emph{Interoperable}: the schema extends GBFS v3
and is published as a JSON Schema, together with an explicit
mapping to the Transport Mode Ontology so that automated
pipelines can consume it without manual configuration. Finally,
\emph{Reusable}: the dataset is released under the Open Database
Licence v1.0 and accompanied by both the audit report and the
source code, which together support an end-to-end reproduction
of the production pipeline.

\paragraph{Versioning}
Every release of the Gold Standard carries a semantic version
number (Major.Minor.Patch) and a Git hash that uniquely
identifies the pipeline that produced it
(Table~\ref{tab:versions}). This is not a cosmetic convention:
versioning is the only way to guarantee that a downstream
analysis can be re-run on the exact dataset that was used in a
given publication, and that ulterior schema evolutions remain
explicit and traceable.

\begin{table}[ht]
\centering
\caption{Versioning of the Gold Standard GBFS France.}
\label{tab:versions}
\small
\begin{tabular}{@{}lll@{}}
\toprule
\textbf{Version} & \textbf{Date} & \textbf{Notes}\\
\midrule
0.1.0 & 2025-04 & Raw ingestion, no purging.\\
0.5.0 & 2025-09 & A1, A4 (partial).\\
0.9.0 & 2026-01 & A2, A3 added; $\sigma_{\max}$ tuned.\\
1.0.0 & 2026-04 & A1--A5 closed; audit report released.\\
\bottomrule
\end{tabular}
\end{table}

\paragraph{Checklist}
A formal reproducibility checklist accompanies the released
dataset: the raw sources are archived with SHA-256 hashes, the
Python environment is pinned in a \texttt{requirements.txt} file,
a Docker container is published for bit-exact reproduction,
random seeds are fixed (\texttt{numpy=42},
\texttt{python=42}), idempotence unit tests are executed in CI,
the audit report is generated automatically on every release,
and code, data and report are linked through a single DOI.

\section{\uppercase{Implications and Limitations}}
\label{secCommunity}

\paragraph{Towards a shared audit protocol}
The audit documented in this paper is intended as a scientific
validation tier, downstream of GBFS publication and upstream of
academic work. Its sustainability, however, depends on collective
appropriation, since no single research team can durably operate
a national infrastructure on its own. Three concrete avenues are
suggested. First, a jointly published \emph{GBFS Audit Best
Practices Handbook} by the francophone micromobility laboratories
(CESI LINEACT, IFSTTAR UMRESTTE, LAET, LVMT, IRSTEA-ITRT) would
align the empirical practices of teams that currently rebuild the
same chain in isolation. Second, a working group placed under
ADEME or \emph{transport.data.gouv.fr} could formalise an
extended GBFS schema and a certification framework, giving the
audit a regulatory anchor. Third, a continuous validation chain
running the pipeline monthly on the national feeds, with
automated publication of the audit report and operator-side
notification on regression, would close the loop between
detection and correction.

\paragraph{Governance}
A data infrastructure survives only if it is backed by a
governance body that decides on schema evolutions, mediates
between contributors and resolves classification conflicts. Star
\cite{Star2010} has shown that the most durable infrastructures
are those whose governance is explicit, transparent and tied to
an identifiable community of practice. The Gold Standard as
released today is, de facto, governed by the authors of the
pipeline: this situation is acceptable as a transitional state,
but it is inadequate in the medium term. Moving to a shared
governance is therefore an institutional challenge as much as a
scientific one, and is a precondition for the protocol to scale
beyond the perimeter of the BikeShare-ICT research programme.

\paragraph{Limitations}
The protocol has, finally, four acknowledged limitations.
\emph{L1 -- National coverage.} The corpus is limited to
metropolitan France, and an extension to overseas territories
and to other European jurisdictions will require revisiting the
thresholds, in particular the macro-regional surface threshold
and the geofilter bounds. \emph{L2 -- Static versus dynamic.}
The current version focuses on \texttt{station\_information};
auditing the \texttt{station\_status} feed (real-time
availability) calls for a distinct protocol that is sketched in
the validation roadmap and developed in a companion paper
\cite{Fosse2026flows}. \emph{L3 -- Free-floating asymmetry.}
The protocol is calibrated for dock-based fleets and free-
floating fleets are preserved in the corpus but not rigorously
audited, since their analysis requires a dedicated framework with
trip-level statistics. \emph{L4 -- Operator coverage.} Some
operators publish partial GBFS feeds or are unavailable during
the audit window, so that the completeness of the inventory is
not guaranteed and must be re-evaluated periodically.

\paragraph{Recommendations to operators}
The audit diagnoses problems but should also turn into a
constructive feedback to the operators that publish GBFS. Four
recommendations follow directly from the taxonomy. First, the
\texttt{vehicle\_type} field of GBFS v3 should be populated
exhaustively and consistently, so that car-sharing or
scooter-sharing services published on the national portal can no
longer be ingested as bike-sharing fleets (A1). Second, the
\texttt{capacity} field should report a physical dock count for
dock-based stations and should be left to \texttt{null} for
virtual stations, rather than filled with placeholders such as
$c = 100$ (A2). Third, free-floating systems should adopt the
GBFS v3 \texttt{vehicle\_status} extension and stop publishing
the per-vehicle anchor positions as \texttt{station\_information}
entries (A3), since the latter introduces a conditional-averaging
bias that no downstream consumer can detect without context.
Fourth, geographical perimeters should be published as polygons
in the upcoming GBFS \texttt{geofencing\_zones} block, which
would let validators apply system-specific topological filters
instead of relying on national-scale heuristics (A4, A5). These
recommendations are deliberately compatible with the current
GBFS specification: they are conventions that the French
operator community could adopt voluntarily, without waiting for
a schema revision.

\paragraph{Downstream use cases}
The Gold Standard is not produced for its own sake but to unlock
a class of analyses that the raw feed would not support. Four
families of downstream studies illustrate the expected return on
investment. The first is \emph{composite-index work} -- the
companion IMD paper \cite{Fosse2026imd} is the immediate
example, but any composite indicator that mixes infrastructure,
multimodality and safety variables at the station level inherits
the same audit. The second is \emph{demand modelling and
forecasting}: graph-based neural networks
\cite{Lin2018BSSGraph} are extremely sensitive to false
\emph{station} nodes introduced by A3, and the Gold Standard
removes that confound without changing the operator-level data
pipeline. The third is \emph{spatial-equity analysis}, where the
hybridisation with INSEE Filosofi opens the kind of
income-versus-access study that has been done sporadically on
single cities but never at the national level. The fourth is
\emph{multimodal integration with GTFS}, since the hybridisation
schema is deliberately compatible with the canonical GTFS audit
toolchain \cite{MobilityData2024Validator}; a single research
team can therefore audit transit and shared-mobility feeds with
the same conceptual framework.

\section{\uppercase{Validation Roadmap}}
\label{secExperiments}

The four limitations of Section~\ref{secCommunity} can each be
turned into a falsification programme. In this section, six
experiments are catalogued, each targeting one or two threats to
validity and specified by (i) a hypothesis, (ii) a protocol that
another team could re-execute, and (iii) an explicit success
criterion, without which any verification attempt would remain
unfalsifiable. The mapping between experiments and limitations
is summarised in Table~\ref{tab:roadmap}.

\begin{table*}[ht]
\centering
\caption{Experiments of the validation roadmap, mapped to the
limitations declared in Section~\ref{secCommunity}.}
\label{tab:roadmap}
\small
\begin{tabular}{@{}llll@{}}
\toprule
\textbf{ID} & \textbf{Title} & \textbf{Targets} &
\textbf{Status}\\
\midrule
E1 & Inter-rater reliability of A1--A5 &
    Auditor effect, A3 detection rule       & Automated pilot\\
E2 & Threshold and buffer sensitivity      &
    $\sigma_{\max}$, $N_{\min}$, radii      & $\sigma_{\max}$ pilot complete\\
E3 & Twelve-month temporal stability       &
    Snapshot dependence (L4), churn         & Running\\
E4 & Dynamic extension to \texttt{station\_status} &
    L2 (static vs.\ dynamic)                & 2-day pilot complete\\
E5 & European generalisation (5 countries) &
    L1 (national coverage)                  & Planned 2027-Q1\\
E6 & Free-floating-native protocol         &
    L3 (free-floating asymmetry)            & Planned 2027-Q2\\
\bottomrule
\end{tabular}
\end{table*}

\paragraph{E1 -- Inter-rater reliability}
The first experiment seeks to defuse the auditor effect, \ie
the suspicion that A1--A5 (and in particular A3) is a set of
rules designed around the very systems it identifies. The
hypothesis is that the taxonomy is sufficiently operational to be
applied by annotators who did not participate in its
construction. The protocol draws a stratified random sample of
500 stations covering all 122 systems, over-sampled for systems
flagged A2--A3, and recruits three independent transport-data
engineers external to the team. Each annotator classifies every
station given only the textual definitions of A1--A5 and the raw
GBFS payload of the station; pairwise Cohen's $\kappa$ and the
three-way Fleiss $\kappa$ are computed, and disagreements are
adjudicated by a fourth blind reviewer. The pipeline is then
re-run with the adjudicated labels as ground truth, and the
precision and recall of the rule-based classifier are reported.
The experiment is considered successful if Fleiss $\kappa \ge
0.75$ across all five classes and the rule-based recall reaches
$\ge 0.90$ on A1--A3.

In the absence of external human annotators at the time of
writing, a complementary automated robustness check was performed
on the 500-station sample. Three independent rule-based
classifiers were implemented, each restricted to a single
orthogonal signal source: a \emph{name-based} classifier seeing
only \texttt{system\_name} and \texttt{station\_name}; a
\emph{capacity-statistics} classifier seeing only the
\texttt{capacity} column; and a \emph{geospatial} classifier
seeing only \texttt{lat} and \texttt{lon}. The pairwise Cohen's
$\kappa$ between the three classifiers is near zero ($\kappa <
0.10$) and Fleiss $\kappa = -0.07$, which would be a failure
condition under the canonical E1 interpretation. It is not: the
three classifiers solve *different* slices of the taxonomy ---
A is the only one that can see A1, B is the only one that can
see A2, C is the only one that can see A4 --- so their
disagreement is the expected signature of \emph{complementary
coverage} rather than annotator inconsistency. The right metric
is therefore the recall of the \emph{union} of the three
classifiers against the rule-based verdict, which reaches
$97.2\,\%$ at a precision of $91.5\,\%$. Only $7$ stations
out of $500$ are flagged by all three classifiers, confirming
that A1--A5 cannot be reduced to a single heuristic and that the
integration step of the purging protocol is itself a substantive
contribution. The full E1 with three external human annotators
is, accordingly, the next priority of the validation roadmap;
the automated coverage check is reported here as a
\emph{robustness} measure, not as a substitute for human
annotation. All three classifiers, the coverage analysis script
and the 500-station blinded sample are released alongside the
data product.

\paragraph{E2 -- Threshold and buffer sensitivity}
The second experiment examines whether the certified set
$\mathcal{G}$ and the downstream ranking remain stable under
reasonable variations of $\sigma_{\max}$, $N_{\min}$, the BAAC
buffer and the BD~TOPO buffer. The Cartesian grid
$(\sigma_{\max}, N_{\min}, r_{\text{BAAC}}, r_{\text{BD~TOPO}})
\in \{2,2.5,3,3.5,4\}
\times \{10,20,30,50\}
\times \{300,500,800\}$\,m
$\times \{200,300,500\}$\,m
is swept; for every grid point, the certified set, the enriched
variables and the IMD ranking are recomputed, and the Jaccard
similarity to the reference run, Kendall's $\tau$ on the ranking
and first- and total-order Sobol indices on $\tau$ are reported.
The success criterion is, on the modal grid (centre $\pm$ one
step), Kendall's $\tau \ge 0.95$ and Jaccard $\ge 0.97$. Any
parameter with a total Sobol index above 0.40 is to be
reclassified from \emph{operational threshold} to
\emph{modelling choice} and documented as such in the schema.

A single-dimension pilot of E2 has been completed on
$\sigma_{\max}$, the parameter for which an a priori suspicion
of fragility had been raised. Sweeping
$\sigma_{\max} \in \{2.0, 2.5, 3.0, 3.5, 4.0\}$ on the certified
corpus yields Jaccard similarities of $0.66, 0.79, 1.00, 1.00,
1.00$ and Kendall's $\tau$ of $0.92, 0.95, 1.00, 1.00, 1.00$
against the reference at $\sigma_{\max} = 3.0$. Most importantly,
the national \textbf{Top-10 is invariant} across the entire grid
($\sigma_{\max}$-induced churn = 0), which is the criterion that
matters for any downstream public-policy reading of the ranking.
The four-dimensional grid sweep over
$(\sigma_{\max}, N_{\min}, r_{\text{BAAC}}, r_{\text{BD~TOPO}})$
requires re-running the hybridisation pipeline at every grid
point and is the next priority.

\paragraph{E3 -- Twelve-month temporal stability}
The third experiment tests whether anomaly verdicts are stable in
time \ie whether the pipeline reclassifies automatically when an
operator publishes a corrected feed but does not introduce
spurious noise when nothing has changed upstream. The audit is
re-executed on the first of every month over twelve consecutive
months, each version is archived with its hash, and for each
system the verdict time series
$v_s(t) \in \{\text{ok}, A1, \ldots, A5\}$ is built, together
with the churn rate
$\rho_s = \frac{1}{T-1}\sum_t \mathbf{1}[v_s(t) \neq v_s(t-1)]$.
Verdict changes are then correlated with announced operator
releases, in order to distinguish an \emph{operator effect}
(legitimate reclassification) from an \emph{audit drift}
(spurious noise produced by the pipeline itself). The experiment
is considered successful if the median $\rho_s$ across the 85
certified systems is below 0.10 and all verdict changes are
traceable to an exogenous event.

\paragraph{E4 -- Dynamic extension}
The fourth experiment hypothesises that a second family of
anomalies, invisible in \texttt{station\_information}, becomes
detectable when the \texttt{station\_status} feed is monitored
over time: counts exceeding declared capacity, zero-availability
persisting for weeks, geo-fence drifts. These are referred to as
A6 candidates. The protocol collects
\texttt{station\_status} snapshots every five minutes for 30 days
on the 85 certified systems, flags stations where
$P(\text{available} > \text{capacity}) > 0.01$, identifies
pathologically static stations
($P(\text{unchanged for 24h}) > 0.40$), and publishes the dynamic
profile as a new field \texttt{audit\_status\_dynamic}. The
experiment is successful if at least three statistically distinct
A6 sub-types are identified across $\ge 5$ systems and the
dynamic audit reclassifies $\ge 5\,\%$ of stations whose static
verdict was \emph{ok}.

A 2-day pilot of E4 has been completed on the 43 systems for
which \texttt{station\_status} snapshots were already collected
(4{,}635 certified stations monitored, 3.0\,M individual
snapshots). Three indicators were computed per station:
$\rho_{\text{overflow}} = P(\text{bikes} > \text{capacity})$,
$\rho_{\text{empty}} = P(\text{bikes} = 0)$ and
$\rho_{\text{full-zero}} =
P(\text{bikes} = 0 \land \text{docks} = 0)$. Applying the
thresholds of $0.01$, $0.40$ and $0.20$ respectively, the pilot
identifies 109 \emph{overflow} cases, 492 \emph{saturated-empty}
cases and 87 \emph{degenerate} cases. The union of the three A6
sub-types covers $599$ stations, \ie $12.9\,\%$ of the monitored
corpus, well above the $5\,\%$ success threshold even on this
short 2-day window. The full 30-day E4 is in progress.

\paragraph{E5 -- European generalisation}
The fifth experiment addresses the question of portability. The
hypothesis is that A1--A5 captures structural features of GBFS
publication and can be transferred to other European
jurisdictions, modulo per-country threshold recalibration; the
alternative hypothesis is that A3, in particular, is an artefact
of the French free-floating market and would not appear
elsewhere, which would invalidate any claim of generality. The
pipeline is therefore applied to the national GBFS feeds of
Spain (MITMA), Italy (MIMS), Germany (Mobilithek), Belgium
(data.gov.be) and the Netherlands (NDOV); E1 is repeated with
three local annotators per country; and $\sigma_{\max}$ and
$N_{\min}$ are recalibrated per country via the E2 sensitivity
grid. The success criterion is that A1, A3 and A4 are detected
in at least four of the five countries with country-level
$\kappa \ge 0.70$, and that recalibrated thresholds remain
within $\pm 20\,\%$ of the French values.

\paragraph{E6 -- Free-floating-native protocol}
The last experiment opens the question of free-floating fleets,
for which the present protocol is acknowledged to be only a first
approximation. The hypothesis is that such fleets warrant a
dedicated audit framework whose anomaly classes are not subsumed
by A1--A5. On the eight free-floating systems of the corpus,
five candidate classes are defined: F1 (geo-fence drift), F2
(dwell-time saturation), F3 (vehicle-id non-persistence), F4
(trip-distance distribution incompatible with declared vehicle
class) and F5 (abandonment outside service area). Each class is
codified as a statistical test on \texttt{free\_bike\_status}
snapshots and on inferred trips, E1-style inter-rater reliability
is run on a 200-vehicle sample, and a complementary Floating
Standard dataset is published. The experiment is successful if
at least three of the five candidate classes survive E1-style
validation ($\kappa \ge 0.70$) and their combined incidence
covers $\ge 60\,\%$ of the free-floating fleet.

\section{\uppercase{Conclusion}}
\label{secConclusion}

This study aimed to address a question that is rarely posed
explicitly in the bike-sharing literature: can an open
standardised feed be ingested directly into quantitative research
pipelines, or does it need an intermediate validation tier? By
auditing the 122 French GBFS systems, we have shown that the
syntactic standardisation of an open feed is not, in itself,
sufficient to turn it into a research artefact. Five recurring
anomaly classes were identified, whose cumulative incidence is
sufficient to invalidate any non-audited comparative ranking, as
emblematically illustrated by Bordeaux's shift from rank 2 to
rank 14 after correcting a single anomaly. Beyond this empirical
diagnosis, three practical lessons can be drawn for the
community. First, the production of a Gold Standard is neither a
standalone product nor a methodological complement to a main
study: it is a scientific contribution in its own right, and
deserves to be cited, valued and reviewed as such. Second, the
sustainability of such an infrastructure depends on collective
appropriation more than on technical performance, and a shared
audit handbook, a working group and a continuous validation chain
would be necessary first steps in that direction. Third, while
the present version of the protocol is technically operational,
the validation roadmap of Section~\ref{secExperiments} is what
would move the Gold Standard from a descriptive artefact to a
certifiable infrastructure, and any failure on one of its six
experiments should be read as an invitation to revise the
protocol in public rather than to defend it in private.

\section*{Data Availability}
The Gold Standard GBFS France data product is released openly
and is deposited on three complementary channels that together
satisfy the FAIR principles for research data.

\emph{Primary deposit (Zenodo).} The dataset is archived on
Zenodo under the concept DOI
\texttt{10.5281/zenodo.XXXXXXX}, with one DOI per
SemVer-tagged release (current version 1.0.0, dated 2026-04).
Each release contains the certified Parquet file, the audit
report and the SHA-256 hashes of the raw GBFS snapshots used as
input. Zenodo guarantees long-term preservation and indexes the
dataset in OpenAIRE.

\emph{Mirror (Recherche Data Gouv).} The dataset is mirrored on
the French national research data repository
\emph{Recherche Data Gouv} (Dataverse instance operated by the
MESR), which provides additional indexing in HAL and the
\emph{Bibliographie Scientifique Ouverte} (BSO).

\emph{Pipeline (GitHub).} The source code of the audit pipeline,
the unit tests (idempotence, completeness) and the schema
definition are hosted at
\url{https://github.com/rohanfosse/bikeshare-data-explorer},
with releases mirrored to Zenodo. The repository ships a
\texttt{CITATION.cff} file and a Docker image
(\texttt{ghcr.io/.../gold-standard:1.0.0}) for bit-exact
reproduction.

\emph{Standards compliance.} The dataset is described by both a
DCAT-AP metadata record \cite{DCATAP2021} (for European
open-data portals) and a Croissant JSON-LD manifest
\cite{Croissant2024} that makes it discoverable by
Google Dataset Search, Hugging Face Datasets and Kaggle.
A Frictionless Data Package descriptor
\cite{FrictionlessData2023} is shipped alongside the Parquet
file, enabling out-of-the-box validation by any frictionless-
aware tool.

\emph{Licences.} The data is released under the Open Database
Licence (ODbL v1.0); the code is released under the MIT
licence. Derivative datasets are required to redistribute the
ODbL terms and to cite the concept DOI.

\emph{Citation.} Authors using the Gold Standard are kindly
asked to cite the concept DOI alongside the present paper. A
\texttt{CITATION.cff} file is provided at the root of the
GitHub repository so that the citation can be exported in
BibTeX, RIS, CSL-JSON or Endnote with a single click.

\section*{Acknowledgments}
The authors thank the producers of open data that made this
research possible: GBFS operators, MobilityData, OpenStreetMap,
Cerema, ONISR, INSEE, FUB and the French GTFS community.

\section*{Declarations}
The authors declare no conflicts of interest.

\begin{comment}
\begin{itemize}
\item Funding
\item Conflict of interest/Competing interests
\item Ethics approval
\item Consent to participate
\item Consent for publication
\item Availability of data and materials
\item Code availability
\item Authors' contributions
\end{itemize}
\end{comment}

\bibliographystyle{unsrt}
\bibliography{references}

\end{document}
